% Secciones principales

\section*{Planteamiento del Proyecto}
El planteamiento o definición correcta del proyecto es esencial para no desviar el objetivo. Aquí se brinda una descripción concreta del proyecto.

\section*{Objetivo}
El objetivo general surge directamente del proyecto y define sus alcances. Es importante tener una redacción clara y precisa.

\section*{Antecedentes}
Este apartado analiza lo que se ha escrito sobre el proyecto, abordando estudios previos y características relevantes a imitar o mejorar.

\section*{Metodología}
La metodología constituye el diseño del proyecto, explicando la forma en que se desarrollará. Todos los pasos deben ser detallados y permitir modificaciones justificadas.

\section*{Cronograma}
Es esencial agregar un cronograma explicando las etapas del proyecto. Lo más conveniente es usar periodos como semanas o meses.

\section*{Recursos}
Listado de las actividades que requieren apoyos como material, equipo, salarios, etc., relacionado con el presupuesto.

\section*{Bibliografía}
Las referencias deben agregarse en el texto para indicar el soporte en otros autores o trabajos.

% Requisitos de formato
\section*{Formato de la Presentación}
\subsection*{A. Requisitos de Formato}
\begin{itemize}
    \item Papel A4 (210 x 297 mm), tinta negra.
    \item Títulos en Arial 14, negrita.
    \item Texto en Arial tamaño 10-11 o Times New Roman 11-12.
\end{itemize}

\subsection*{B. Márgenes}
\begin{itemize}
    \item Izquierdo: 3.0 cm
    \item Superior: 3.0 cm
    \item Derecho: 1.0 cm
    \item Inferior: 2.0 cm
    \item Encuadernación: 1.0 cm
\end{itemize}

\subsection*{C. Paginación}
Números arábigos en la parte inferior derecha, desde la introducción hasta el final del documento.

\subsection*{D. Espaciado}
Interlineado de 1.5 espacios.

\subsection*{E. Alineación de Texto}
\begin{itemize}
    \item Sangría normal en la primera línea de cada párrafo.
    \item Alineación justificada.
\end{itemize}
